%% kapa: exam=1 round=29 year=2018 period=2 form=A subject=appraisal_law subject_ko=감정평가관계법규 seq=30 number=30
%% source: Q-Net 감정평가사 기출문제 (cst00309) artlSeq=5202808 file=extracted/1008/5202808/2018년 제29회 감정평가사 1차시험 기출문제/2018년 제29회 감정평가사 1차시험 기출문제/2교시 A형.pdf
%% conversion: kordoc 4.24.51 -> tools/gen1.py (auto); review=proofread-sonnet; flags=none
\begin{question}{30}
공간정보의 구축 및 관리 등에 관한 법령상 지적공부에 관한 설명으로 옳지 않은 것은?
\begin{choices}
  \item 지적공부를 정보처리시스템을 통하여 기록ㆍ저장한 경우 그 지적공부는 지적 정보관리체계에 영구히 보존되어야 한다.
  \item 정보처리시스템을 통하여 기록ㆍ저장된 공유지연명부를 열람하려는 경우에는 특별자치시장, 시장ㆍ군수 또는 구청장이나 읍ㆍ면ㆍ동의 장에게 신청할 수 있다.
  \item 지방자치단체장이 지적전산자료를 신청하는 경우에는 지적전산자료의 이용에 관하여 미리 관계 중앙행정기관의 심사를 받아야 한다.
  \item 정보처리시스템을 통하여 기록ㆍ저장된 지적공부 사항 중 소유자에 관한 사항을 복구할 때에는 부동산등기부나 법원의 확정판결에 따라야 한다.
  \item 시ㆍ군ㆍ구 단위의 지적전산자료를 이용하거나 활용하려는 자는 지적소관청에 지적전산자료를 신청하여야 한다.
\end{choices}
\end{question}
\answer{3}
